%% Dans ce fichier, les �ventuels accents sont cod�s
%% en ISO-8859-15 (aka latin9)
%% La ligne suivante donne trois exemples 
%% e accent aigu : �, e accent grave : �, e-dans-l'o : �
%%%%%ENTREES%%%%
\begin{tabarg}
\multicolumn{3}{l}{\titretab {Entr\'ees}}\\
\hline
$U$ & Vecteur, \texttt{float} & 1�re composante du vent horizontal (m/s) \\
$V$ & Vecteur, \texttt{float} & 2�me composante du vent horizontal (m/s) \\
$W$ & Vecteur, \texttt{float} & Composante du vent vertical (m/s) \\
$N_{s}$ & Coeff., \texttt{int} & Nombre de sous-p�riodes par p�riode\\
$p$ & Coeff., \texttt{float} & pourcentage de donn�es en de�a duquel le calcul ne sera pas effectu� (valable pour p�riodes et sous-p�riodes)\\
$time_{out}$ & Vecteur, \texttt{double} & Base de temps d'arriv�e (secondes) dont le time\_bnds d�finit les p�riodes de calcul.\\
\end{tabarg}

\medskip
%%%%%SORTIES%%%%
\begin{tabarg}
\multicolumn{3}{l}{\titretab {Sorties}}\\
\hline
$I_{stat}$ & Vecteur, \texttt{float} & Indice de stationnarit� (sans unit�)\\
\end{tabarg}
